% ======================================================================
% Chapter 6: Radiative Corrections: Introduction
% (book pages 175-210; PDF pages 198-233)
% ======================================================================
\chapter{辐射修正：引言}

既然我们已经获得了进行 QED 计算的一些经验，让我们转向一些更复杂的问题。第~5 章只处理了树图过程，即不含圈的图。但所有这些过程都从确实含有圈的图中得到高阶贡献，称为\emph{辐射修正}。QED 中辐射修正的另一个来源是\emph{轫致辐射}，即反应过程中额外末态光子的发射。在本章中，我们将研究这两类辐射修正，并且将会发现，只包含其中一种而不包含另一种是不自洽的。

在整个本章中，为了在最简单的情形下说明这些概念，我们将考虑电子在另一个非常重的粒子上的散射过程。我们在第~5.4 节和习题~5.1 中已经在树图级别分析了这一过程。在微扰论的下一阶，我们遇到以下四个图：
\begin{equation}
  \text{(vertex correction)} + \text{(two external leg corrections)}
    + \text{(vacuum polarization)} .
  \label{eq:6.1}
\end{equation}
截面的 $\alpha$ 阶修正来自这些图与树图之间的干涉项。还有六个涉及圈中重粒子的一圈图，但在该粒子远重于电子的极限下可以忽略它们，因为质量出现在传播子的分母中。（物理上，重粒子在碰撞中加速较小，因此辐射也较少。）

在 \eqref{eq:6.1} 的四个图中，第一个（称为\emph{顶角修正}）最为复杂，并给出最多样的新效应。例如，它导致电子的反常磁矩，我们将在第~6.3 节中计算它。

\eqref{eq:6.1} 中接下来的两个图是\emph{外腿修正}。在本章中我们将忽略它们，因为它们没有按我们计算 $S$ 矩阵元的公式式~(4.90) 的要求被截断。当我们将在第~7.2 节证明那个公式时，会更详细地讨论这些图。

\eqref{eq:6.1} 中最后一个图称为\emph{真空极化}。由于它比其他图需要更多的计算技术，我们要到第~7.5 节才会计算这个图。

对这些修正的研究会因为这些修正本身无定义而变得复杂。\eqref{eq:6.1} 中的每个图都涉及对未确定的圈动量的积分，并且在每种情形下，积分在 $k\to\infty$ 或\emph{紫外}区域都是发散的。幸运的是，这些积分的无穷大部分总会从诸如截面之类的可观测量表达式中消去。

\eqref{eq:6.1} 中前三个图还包含\emph{红外发散}：来自圈动量积分 $k\to 0$ 一端的无穷大。我们将在第~6.4 节看到，当我们同时包括以下\emph{轫致辐射}图时，这些发散会被抵消：
\begin{equation}
  \text{photon emission by initial electron} + \text{photon emission by
    final electron} .
  \label{eq:6.2}
\end{equation}
这些图在辐射光子能量趋于零的极限下是发散的。在这个极限下，任何物理探测器都无法观测到该光子，因此把产生这些低能光子的截面加到无辐射散射截面上是有意义的。因此，轫致辐射图是辐射修正的基本组成部分，无论在此处还是在任何其他 QED 过程中都是如此。

本章的主要目标是理解低能光子的轫致辐射、顶角修正图，以及这两类辐射修正之间红外发散的抵消。

\section{软轫致辐射}

让我们从分析轫致辐射过程开始对辐射修正的研究。在本节中，我们首先对电子突然加速时低频轫致辐射的强度做经典计算。然后我们将在量子场论中计算一个密切相关的量：由 \eqref{eq:6.2} 中诸图给出的发射一个非常软的光子的截面。我们想要理解经典结果如何作为量子结果的极限情形而出现。

\subsection{经典计算}

设一个经典电子在时刻 $t=0$、位置 $\mathbf{x}=0$ 处受到一次突然的冲击，使其 4-动量从 $p$ 变为 $p'$。（无穷突然的动量变化当然是一种不现实的理想化。然而，加速期间轨迹的精确形式并不影响低频辐射。我们的计算对频率小于散射时间倒数的辐射有效。）

我们可以写下这个电子的电流，并把该电流视为麦克斯韦方程组的源，从而求得辐射场。

这样一个粒子的电流密度是什么？对于静止在 $\mathbf{x}=0$ 处的带电粒子，电流为
\begin{equation}
  j^\mu(\mathbf{x}) = e\,\delta^{(3)}(\mathbf{x})\,
    \begin{pmatrix}1\\ 0\end{pmatrix}^{\mu}
  = e\int dt\ \begin{pmatrix}1\\ 0\end{pmatrix}^{\mu}\,
    \delta^{(3)}\bigl( \mathbf{x} - \mathbf{y}(t) \bigr)
  \qquad\text{with } y^\mu(t) = (t, 0)^\mu .
\end{equation}
由此我们可以猜出任意轨迹 $y^\mu(\tau)$ 的电流：
\begin{equation}
  j^\mu(x) = e\int d\tau\ \frac{dy^\mu}{d\tau}\,
    \delta^{(4)}\bigl( x - y(\tau) \bigr).
  \label{eq:6.3}
\end{equation}
注意这个表达式与曲线 $y^\mu(\tau)$ 被参数化的精确方式无关：把变量从 $\tau$ 换到 $a(\tau)$ 会在积分测度中给出一个因子 $d\tau/da$，它通过链式法则与 $dy^\mu/d\tau$ 结合给出 $dy^\mu/da$。我们还可以从 \eqref{eq:6.3} 证明该电流自动守恒：对任意在无穷远处衰减的``试探函数'' $f(x)$，我们有
\begin{equation}
  \begin{aligned}
    \int d^4x\ f(x)\,\partial_\mu j^\mu(x)
    &= \int d^4x\ f(x)\,e\int d\tau\
      \frac{dy^\mu}{d\tau}\ \frac{\partial}{\partial x^\mu}
      \delta^{(4)}(x-y(\tau))\\
    &= -e\int d\tau\ \frac{dy^\mu}{d\tau}\ \frac{\partial f}{\partial x^\mu}
      \Big|_{x=y(\tau)}
    = -e\int d\tau\ \frac{d}{d\tau}f\bigl( y(\tau) \bigr)\\
    &= -e\Bigl[ f\bigl( y(\tau) \bigr) \Bigr]_{-\infty}^{\infty} = 0 .
  \end{aligned}
\end{equation}
对我们的过程，轨迹为
\begin{equation}
  y^\mu(\tau) = \begin{cases}
    (p^\mu/m)\,\tau & \text{for } \tau < 0;\\
    (p'^{\mu}/m)\,\tau & \text{for } \tau > 0 .
  \end{cases}
\end{equation}
于是电流可以写成
\begin{equation}
  j^\mu(x) = e\int_0^{\infty} d\tau\ \frac{p'^{\mu}}{m}\,
    \delta^{(4)}\Bigl( x - \frac{p'}{m}\tau \Bigr)
    + e\int_{-\infty}^{0} d\tau\ \frac{p^{\mu}}{m}\,
      \delta^{(4)}\Bigl( x - \frac{p}{m}\tau \Bigr).
\end{equation}
稍后我们将需要知道这个函数的 Fourier 变换。插入因子 $e^{-\epsilon\tau}$ 和 $e^{\epsilon\tau}$ 以使积分收敛，我们有
\begin{equation}
  \begin{aligned}
    \tilde{j}^\mu(k) &= \int d^4x\ e^{ik\cdot x}\,j^\mu(x)\\
    &= e\int_0^{\infty} d\tau\ \frac{p'^{\mu}}{m}\,
      e^{i(k\cdot p'/m+i\epsilon)\tau}
      + e\int_{-\infty}^{0} d\tau\ \frac{p^{\mu}}{m}\,
        e^{i(k\cdot p/m-i\epsilon)\tau}\\
    &= ie\Bigl( \frac{p'^{\mu}}{k\cdot p' + i\epsilon}
      - \frac{p^{\mu}}{k\cdot p - i\epsilon} \Bigr).
  \end{aligned}
  \label{eq:6.4}
\end{equation}
我们现在准备好求解麦克斯韦方程组。在 Lorentz 规范（$\partial_\mu A^\mu = 0$）中，我们必须求解
\begin{equation}
  \partial^2 A^\mu = -j^\mu,
\end{equation}
或换言之在 Fourier 空间中，
\begin{equation}
  \tilde{A}^\mu(k) = -\frac{1}{k^2}\tilde{j}^\mu(k) .
\end{equation}
将 \eqref{eq:6.4} 代入，我们得到矢势的公式：
\begin{equation}
  A^\mu(x) = \int \frac{d^4k}{(2\pi)^4}\ e^{-ik\cdot x}\
    \frac{-ie}{k^2}\ \Bigl( \frac{p'^{\mu}}{k\cdot p' + i\epsilon}
      - \frac{p^{\mu}}{k\cdot p - i\epsilon} \Bigr).
  \label{eq:6.5}
\end{equation}
$k^0$ 积分可以作为复平面中的围道积分来进行。极点的位置如下：在 $k^0 = \pm|\mathbf{k}|$（来自 $1/k^2$ 因子）以及在 $k\cdot p = 0$ 和 $k\cdot p' = 0$ 处有极点。我们把 $k^0 = \pm|\mathbf{k}|$ 处的极点放在实轴下方，以便（我们很快就会确认）辐射场满足推迟边界条件。

对 $t < 0$，我们向上闭合围道，拾取 $k\cdot p = 0$ 处的极点，即 $k^0 = \mathbf{k}\cdot\mathbf{p}/p^0$。结果为
\begin{equation}
  A^\mu(\mathbf{x}, t) \sim \int \frac{d^3k}{(2\pi)^3}\ e^{i\mathbf{k}\cdot\mathbf{x}}
    e^{-i(\mathbf{k}\cdot\mathbf{p}/p^0)t}\ (2\pi i)\Bigl( \frac{1}{2p^0}
      \cdot \frac{p^\mu}{-k^2}\cdots \Bigr).
\end{equation}
在粒子最初静止的参考系中，其动量矢量为 $p^\mu = (p^0, 0)^\mu$，矢势化为
\begin{equation}
  A^0(\mathbf{x}) = \int \frac{d^3k}{(2\pi)^3}\ e^{i\mathbf{k}\cdot\mathbf{x}}\
    \frac{e}{\mathbf{k}^2} = \frac{e}{4\pi|\mathbf{x}|} .
\end{equation}
这正是未被加速电荷的库仑势。正如我们所预期，在粒子被散射之前不存在辐射场。

散射之后（$t > 0$），我们向下闭合围道，拾取实轴下方的三个极点。$k^0 = \mathbf{k}\cdot\mathbf{p}'/p'^0$ 处的极点给出出射粒子的库仑势。因此，另外两个极点完全负责辐射场。它们的贡献给出
\begin{equation}
  A^\mu_{\mathrm{rad}}(\mathbf{x}, t) =
    \int \frac{d^3k}{(2\pi)^3}\ e^{ik\cdot x}\ \mathcal{A}^\mu(k),
  \label{eq:6.6}
\end{equation}
其中动量空间振幅 $\mathcal{A}^\mu(k)$ 由下式给出
\begin{equation}
  \mathcal{A}^\mu(k) = -ie\ \Bigl( \frac{p'^{\mu}}{k\cdot p'}
    - \frac{p^{\mu}}{k\cdot p} \Bigr).
  \label{eq:6.7}
\end{equation}
（条件 $k^0 = |\mathbf{k}|$ 在此处及本计算的其余部分都是隐含的。）

为了计算辐射的能量，我们必须求出电场和磁场。最简单的方式是把 $\mathbf{E}$ 和 $\mathbf{B}$ 写成复 Fourier 积分的实部，正如我们对 $A^\mu$ 所做的那样：
\begin{equation}
  \mathbf{E}(\mathbf{x},t) = \mathrm{Re}\int \frac{d^3k}{(2\pi)^3}\,
    \bm{\mathcal{E}}(k)\,e^{ik\cdot x},\qquad
  \mathbf{B}(\mathbf{x},t) = \mathrm{Re}\int \frac{d^3k}{(2\pi)^3}\,
    \bm{\mathcal{B}}(k)\,e^{ik\cdot x},
\end{equation}
其中
\begin{equation}
  \bm{\mathcal{E}}(k) = ik^0\bm{\mathcal{A}}(k)
    - i\mathbf{k}\mathcal{A}^0(k),\qquad
  \bm{\mathcal{B}}(k) = i\mathbf{k}\times\bm{\mathcal{A}}(k) .
\end{equation}
辐射的能量为
\begin{equation}
  \text{Energy} = \int dt \int d^2x\ \hat{\mathbf{n}}\cdot
    \bigl( \mathbf{E}\times\mathbf{B} \bigr) .
\end{equation}
完成积分（细节见正文），我们得到
\begin{equation}
  \text{Energy} = \int \frac{d^3k}{(2\pi)^3}\ |\mathbf{k}|\,
    \Bigl( \frac{e^2}{4\pi} \Bigr)\ \mathcal{I}(\mathbf{v}, \mathbf{v}')
  = \frac{\alpha}{4\pi^2}\int_0^{\omega_{\mathrm{max}}} d\omega\,
    \omega \int d\Omega\ \mathcal{I}(\mathbf{v}, \mathbf{v}'),
  \label{eq:6.12}
\end{equation}
其中 $\mathcal{I}(\mathbf{v},\mathbf{v}')$ 是初、末态电子速度的已知函数。

现在让我们计算函数 $\mathcal{I}(\mathbf{v},\mathbf{v}')$。对非相对论运动，$|\mathbf{v}| \ll 1$，我们可以近似
\begin{equation}
  \mathcal{I}(\mathbf{v},\mathbf{v}') \approx
    \frac{1}{4\pi}\ \Bigl( \frac{\mathbf{v}'}{1 - \hat{\mathbf{k}}\cdot\mathbf{v}'}
      - \frac{\mathbf{v}}{1 - \hat{\mathbf{k}}\cdot\mathbf{v}} \Bigr)^2 .
  \label{eq:6.13}
\end{equation}
辐射强度的经典结果通过对该表达式积分得到。对角度积分为
\begin{equation}
  \int \frac{d\Omega}{4\pi}\ \Bigl( \frac{\mathbf{v}}{1-\hat{\mathbf{k}}\cdot\mathbf{v}}
    \Bigr)^2 = \frac{v^2}{\gamma^2(1-v^2)}
    = v^2\,\frac{1-v^2}{(1-v^2)^2}\ \cdots ,
\end{equation}
交叉项可以利用下式求值
\begin{equation}
  \int \frac{d\Omega}{4\pi}\ \frac{1}{(1-\hat{\mathbf{k}}\cdot\mathbf{v})
    (1-\hat{\mathbf{k}}\cdot\mathbf{v}')}
  = \frac{1}{|\mathbf{v}-\mathbf{v}'|}\ \mathrm{arctanh}\,
    \frac{|\mathbf{v}-\mathbf{v}'|}{1 - \mathbf{v}\cdot\mathbf{v}'}
    \cdot\ \cdots .
  \label{eq:6.16}
\end{equation}
这些积分很容易完成，我们得到
\begin{equation}
  \mathcal{I}(\mathbf{v},\mathbf{v}') \approx
    \frac{1}{2}\Bigl( \frac{\mathbf{v}'^2}{1-v'^2} + \frac{\mathbf{v}^2}{1-v^2}
      \Bigr) - \frac{1}{2}\ \cdots ,
  \label{eq:6.17}
\end{equation}
其中 $q^2 = (p'-p)^2$。

总之，我们发现低频辐射能量由下式给出
\begin{equation}
  \text{Energy} = \frac{\alpha}{\pi}\int_0^{\omega_{\mathrm{max}}}
    d\omega\ \mathcal{I}(\mathbf{v},\mathbf{v}')
  \xrightarrow{q^2 \gg m^2}
  \frac{\alpha}{\pi}\ \int_0^{\omega_{\mathrm{max}}} d\omega\
    \log\Bigl( \frac{q^2}{m^2} \Bigr).
  \label{eq:6.18}
\end{equation}
如果这份能量由光子构成，每个光子贡献能量 $k$。我们于是预期
\begin{equation}
  \text{Number of photons} = \frac{\alpha}{\pi}\ \int_0^{\omega_{\mathrm{max}}}
    \frac{dk}{k}\ \mathcal{I}(\mathbf{v},\mathbf{v}') .
  \label{eq:6.19}
\end{equation}
我们希望量子力学计算将证实这一结果。

\subsection{量子计算}

现在考虑电子散射期间辐射一个光子的量子力学过程：
\begin{equation}
  e^{-}(p) \to e^{-}(p') + \gamma(k) .
\end{equation}
设 $\mathcal{M}_0$ 表示振幅中来自电子与外场相互作用的部分。则整个过程的振幅为
\begin{equation}
  i\mathcal{M} = -ie\,\bar{u}(p')\ \Bigl[
    \gamma^\mu\epsilon^{*}_\mu(k)\ \frac{i(\slashed{p}'+\slashed{k}+m)}
      {(p'+k)^2-m^2}\ \mathcal{M}_0(p'+k, p)
    + \mathcal{M}_0(p', p-k)\ \frac{i(\slashed{p}-\slashed{k}+m)}
      {(p-k)^2-m^2}\ \gamma^\mu\epsilon^{*}_\mu(k) \Bigr] u(p) .
  \label{eq:6.20}
\end{equation}
由于我们感兴趣的是与经典极限的联系，假设辐射的光子是软的：$|\mathbf{k}| \ll |\mathbf{p}'-\mathbf{p}|$。于是我们可以近似
\begin{equation}
  \mathcal{M}_0(p', p-k) \approx \mathcal{M}_0(p'+k, p)
    \approx \mathcal{M}_0(p', p),
  \label{eq:6.21}
\end{equation}
并且可以忽略传播子分子中的 $\slashed{k}$。分子可以用一些 Dirac 代数进一步简化。在第一项中我们有
\begin{equation}
  (\slashed{p}'+m)\slashed{\epsilon}^{*}u(p) = \bigl[
    2p'\cdot\epsilon^{*} + \slashed{\epsilon}^{*}(\slashed{p}'-m) \bigr]u(p)
    = 2p'\cdot\epsilon^{*}\, u(p) .
\end{equation}
类似地，在第二项中，
\begin{equation}
  \bar{u}(p')\slashed{\epsilon}^{*}(\slashed{p}+m) = \bar{u}(p')\,2p\cdot\epsilon^{*}.
\end{equation}
传播子的分母也化简：
\begin{equation}
  (p-k)^2 - m^2 = -2p\cdot k;\qquad
  (p'+k)^2 - m^2 = 2p'\cdot k .
\end{equation}
因此在软光子近似下，振幅变为
\begin{equation}
  i\mathcal{M} = \bar{u}(p')\mathcal{M}_0(p',p)u(p)\,
    (-ie)\ \Bigl( \frac{p'\cdot\epsilon^{*}}{p'\cdot k}
      - \frac{p\cdot\epsilon^{*}}{p\cdot k} \Bigr).
  \label{eq:6.22}
\end{equation}
这正是弹性散射（无轫致辐射）的振幅乘以一个光子发射因子（方括号内）。

我们过程的截面也很容易用弹性截面来表达；只需为光子变量 $k$ 插入一个额外的相空间积分。对两个光子极化态求和，我们有
\begin{equation}
  \begin{aligned}
    d\sigma(p\to p' + \gamma) &= d\sigma(p\to p')\cdot
      \int \frac{d^3k}{(2\pi)^3}\ \frac{1}{2k}\ \frac{\alpha}{2\pi}\\
    &\quad \times \sum_{\lambda=1,2}\ \Bigl|
      \frac{p'\cdot\epsilon^{(\lambda)}}{p'\cdot k}
      - \frac{p\cdot\epsilon^{(\lambda)}}{p\cdot k} \Bigr|^2 .
  \end{aligned}
  \label{eq:6.23}
\end{equation}
因此在电子从 $p$ 散射到 $p'$ 的条件下辐射一个动量为 $k$ 的光子的微分概率为
\begin{equation}
  d(\mathrm{prob}) = \frac{\alpha}{(2\pi)^2}\ \frac{d^3k}{k}\
    \mathcal{I}(\mathbf{v},\mathbf{v}') .
  \label{eq:6.24}
\end{equation}
这看起来很熟悉；如果我们乘以光子能量 $k$ 来计算期望的辐射能量，就重新得到经典表达式 \eqref{eq:6.12}。

但是有一个问题。式 \eqref{eq:6.24} 不是辐射光子期望数的表达式，而是辐射单个光子的概率的表达式。如果我们对光子动量积分，问题会变得更糟。与 \eqref{eq:6.16} 一样，我们只能积分到软光子近似失效的能量；对这个能量的合理估计是 $|\mathbf{q}| = |\mathbf{p}-\mathbf{p}'|$。因此积分是
\begin{equation}
  \text{Total probability} \approx \frac{\alpha}{\pi}\int \frac{dk}{k}\
    \mathcal{I}(\mathbf{v}, \mathbf{v}') .
  \label{eq:6.25}
\end{equation}
由于 $\mathcal{I}(\mathbf{v},\mathbf{v}')$ 与 $k$ 无关，积分在其下限（我们所有近似在此都很合理）发散。换言之，辐射一个非常软的光子的总概率是无穷大。这就是 QED 微扰论中著名的\emph{红外发散}问题。

我们可以人为地让 \eqref{eq:6.25} 中的积分有定义，办法是假想光子具有一个非常小的质量 $\mu$。这个质量于是为积分提供一个下限截断，使我们能够把本节的结果写成
\begin{equation}
  \begin{aligned}
    d\sigma(p\to p'+\gamma(k)) &= d\sigma(p\to p')\cdot
      \frac{\alpha}{2\pi}\log\Bigl( \frac{-q^2}{\mu^2} \Bigr)
      \mathcal{I}(\mathbf{v},\mathbf{v}')\\
    &\xrightarrow{-q^2\gg m^2}
    d\sigma(p\to p')\cdot \frac{\alpha}{2\pi}
      \log\Bigl( \frac{-q^2}{\mu^2} \Bigr)\log\Bigl( \frac{-q^2}{m^2} \Bigr).
  \end{aligned}
  \label{eq:6.26}
\end{equation}
这一结果对 $q^2$ 的依赖称为\emph{苏达科夫双对数}，它是有物理意义的，并将在第~6.4 节再次出现。然而，对 $\mu$ 的依赖提出了一个我们必须解决的问题。不难猜出，这个问题的解决将涉及把 \eqref{eq:6.24} 重新解释为辐射光子的期望数，而不是辐射单个光子的概率。我们将在第~6.4 节和第~6.5 节看到这种重新解释如何从费曼图得出。然而，为了准备该讨论，我们需要加深对无辐射散射振幅的理解。

\section{电子顶角函数：形式结构}

在简要讨论了由光子发射（轫致辐射）引起的 QED 辐射修正之后，现在让我们研究由于额外虚光子的存在而对电子散射产生的修正：
\begin{equation}
  \text{(one-loop vertex diagram)} .
  \label{eq:6.27}
\end{equation}
这将是我们第一次接触含有圈的费曼图。这类图在量子场论中引起重大而深刻的复杂性。

计算这个图的结果将相当复杂，因此预先思考我们预期这一修正采取什么形式以及如何解释其各种可能的项是有用的。在本节中，我们将考虑顶角修正图的一般性质。我们将看到，Lorentz 不变性、QED 的离散对称性和 Ward 恒等式的基本要求强烈约束着顶角的形式。

那么，考虑这样一类图
\begin{equation}
  \text{(electron-photon vertex)} = -ie\Gamma^\mu(p', p),
\end{equation}
其中灰色圆圈表示最低阶电子-光子顶角与所有截断圈修正之和。我们将这个顶角图之和称为 $-ie\Gamma^\mu(p',p)$。那么，根据我们计算 $S$ 矩阵元的主公式 \eqref{eq:4.103}，电子在重靶上散射的振幅为
\begin{equation}
  i\mathcal{M} = ie^2\ \bigl( \bar{u}(p')\Gamma^\mu(p',p)u(p) \bigr)\
    \frac{1}{q^2}\ \bigl( \bar{u}(k')\gamma_\mu u(k) \bigr).
  \label{eq:6.28}
\end{equation}
更一般地，函数 $\Gamma^\mu(p',p)$ 出现在电子在外部电磁场上散射的 $S$ 矩阵元中。与习题~4.4 一样，向 QED 的哈密顿量加上相互作用
\begin{equation}
  \Delta H_{\mathrm{int}} = \int d^3x\ e\,j^\mu(x)\,A^{cl}_\mu(x),
  \label{eq:6.29}
\end{equation}
其中 $j^\mu(x) = \bar{\psi}(x)\gamma^\mu\psi(x)$ 是电磁流，$A^{cl}_\mu$ 是固定的经典势。在微扰论的最低阶，在该场上的散射的 $S$ 矩阵元为
\begin{equation}
  i\mathcal{M}(2\pi)\delta(p^{0'}-p^0) = -ie\,\bar{u}(p')\gamma^\mu u(p)\,
    \tilde{A}^{cl}_\mu(p'-p),
\end{equation}
其中 $\tilde{A}^{cl}_\mu(q)$ 是 $A^{cl}_\mu(x)$ 的 Fourier 变换。顶角修正把这个表达式修改为
\begin{equation}
  i\mathcal{M}(2\pi)\delta(p^{0'}-p^0) = -ie\,\bar{u}(p')\Gamma^\mu(p',p)u(p)\,
    \tilde{A}^{cl}_\mu(p'-p) .
  \label{eq:6.30}
\end{equation}
在写出 \eqref{eq:6.28} 和 \eqref{eq:6.30} 时，我们有意省略了真空极化图的贡献，例如 \eqref{eq:6.1} 中的第四个图。省略的理由是，这些图应被视为光子传播子的修正，而不是电子-光子顶角的修正。

由于 $\Gamma^\mu(p',p)$ 是 Dirac 空间中的 $4\times4$ 矩阵，我们可以把它表示成第~3.4 节的十六个基矩阵的线性组合。Lorentz 不变性以及离散对称性 $C$、$P$ 和 $T$ 的要求限制了允许的项。人们发现，与这些对称性相容的最一般形式为
\begin{equation}
  \Gamma^\mu(p',p) = \gamma^\mu\,F_1(q^2)
    + \frac{i\sigma^{\mu\nu}q_\nu}{2m}\,F_2(q^2),
  \label{eq:6.33}
\end{equation}
其中 $q = p'-p$，$F_1$ 和 $F_2$ 是 $q^2$ 的标量函数（$F_1(0) = 1$），而 $F_2$ 项是\emph{反常磁矩}部分。为看清这一点，利用 Gordon 恒等式 \eqref{eq:6.32}，
\begin{equation}
  \bar{u}(p')\gamma^\mu u(p) = \bar{u}(p')\Bigl(
    \frac{p'^{\mu}+p^{\mu}}{2m} + \frac{i\sigma^{\mu\nu}q_\nu}{2m} \Bigr)u(p),
  \label{eq:6.32}
\end{equation}
它表明 \eqref{eq:6.33} 中的 $\sigma^{\mu\nu}$ 项贡献一个额外的磁矩。$\Gamma^\mu$ 必须满足 Ward 恒等式的约束
\begin{equation}
  q_\mu\Gamma^\mu(p',p) = 0
\end{equation}
意味着 $F_2$ 完全不受 Ward 恒等式约束，而 $F_1$ 满足 $F_1(0) = 1$（在重整化之后）。

\section{电子顶角函数：计算}

既然我们知道答案应取什么形式（式~\eqref{eq:6.33}），我们准备好计算电子顶角函数的一圈贡献。按如下方式在图上指定动量：入射电子动量为 $p$，出射电子动量为 $p'$，虚光子动量为 $k-p$（于是 $q = p'-p$ 是动量转移）。应用费曼规则，我们发现在 $\alpha$ 阶 $\Gamma^\mu = \gamma^\mu + \delta\Gamma^\mu$，其中
\begin{equation}
  \begin{aligned}
    \bar{u}(p')\,\delta\Gamma^\mu(p',p)\,u(p)
    &= \int \frac{d^4k}{(2\pi)^4}\
      \frac{-ig_{\nu\rho}}{(k-p)^2+i\epsilon}\ \
      \bar{u}(p')(-ie\gamma^\nu)\
      \frac{i(\slashed{k}'+m)}{k'^2-m^2+i\epsilon}\
      \frac{i(\slashed{k}+m)}{k^2-m^2+i\epsilon}\
      (-ie\gamma^\rho)\,u(p)\\
    &= 2e^2\int \frac{d^4k}{(2\pi)^4}\
      \frac{\bar{u}(p')\bigl[ \gamma^\mu k^2 + m^2\gamma^\mu
        - 2m(k+k')^\mu \bigr]u(p)}
      {((k-p)^2+i\epsilon)(k'^2-m^2+i\epsilon)(k^2-m^2+i\epsilon)} .
  \end{aligned}
  \label{eq:6.38}
\end{equation}
在第二行我们使用了缩并恒等式 $\gamma^\nu\gamma^\mu\gamma_\nu = -2\gamma^\mu$。注意分母中的 $+i\epsilon$ 项不能丢弃；它们对于正确计算圈动量积分是必要的。

这个积分看起来不可能，事实上也确实不容易。计算这类积分需要另一项计算技术，称为\emph{费曼参数方法}（尽管一个非常类似的方法更早由 Schwinger 引入）。

\subsection{费曼参数}

这个方法的目的是把 \eqref{eq:6.38} 的三个分母因子压缩成 $k$ 的三次幂单个二次多项式。然后我们可以把 $k$ 平移一个常数以在这个多项式中配平方，并毫无困难地完成剩下的球对称积分。代价将是引入需要被积掉的辅助参数。

最容易的做法是从分母有两个因子的更简单情形开始。我们于是使用恒等式
\begin{equation}
  \frac{1}{AB} = \int_0^1 dx\ \frac{1}{[xA+(1-x)B]^2}
    = \int dx\,dy\ \delta(x+y-1)\ \frac{1}{[xA+yB]^2} .
  \label{eq:6.39}
\end{equation}
它的一个使用示例可能如下：
\begin{equation}
  \frac{1}{(k-p)^2(k^2-m^2)}
  = \int dx\,dy\ \delta(x+y-1)\
    \frac{1}{[k^2 - 2xk\cdot p + xp^2 - ym^2]^2} .
\end{equation}
如果我们现在令 $\ell = k - xp$，我们看到分母只依赖于 $\ell^2$。现在对 $d^4k$ 积分会容易得多，因为 $d^4k = d^4\ell$ 且被积函数关于 $\ell$ 是球对称的。使这一变换成为可能的变量 $x$ 和 $y$ 称为\emph{费曼参数}。

我们的积分 \eqref{eq:6.38} 涉及一个有三个因子的分母，所以我们需要一个稍好一些的恒等式。对 $B$ 微分 \eqref{eq:6.39}，容易证明
\begin{equation}
  \frac{1}{AB^n} = \int dx\,dy\ \delta(x+y-1)\
    \frac{n\,y^{n-1}}{[xA+yB]^{n+1}} .
  \label{eq:6.40}
\end{equation}
但这仍然不够好。我们需要的公式是
\begin{equation}
  \frac{1}{A_1A_2\cdots A_n}
  = \int_0^1 dx_1\cdots dx_n\ \delta\Bigl( \textstyle\sum x_i - 1 \Bigr)\
    \frac{(n-1)!}{[x_1A_1 + x_2A_2 + \cdots + x_nA_n]^n} .
  \label{eq:6.41}
\end{equation}
这个恒等式的证明用归纳法。$n=2$ 的情形就是式~\eqref{eq:6.39}；归纳步骤并不困难，并用到 \eqref{eq:6.40}。

通过对 \eqref{eq:6.41} 反复微分，你可以导出更一般的恒等式
\begin{equation}
  \frac{1}{A_1^{m_1}A_2^{m_2}\cdots A_n^{m_n}}
  = \int_0^1 dx_1\cdots dx_n\ \delta\Bigl( \textstyle\sum x_i - 1 \Bigr)\
    \frac{\Gamma(m_1+\cdots+m_n)}{\Gamma(m_1)\cdots\Gamma(m_n)}\
    \frac{x_1^{m_1-1}\cdots x_n^{m_n-1}}
      {[x_1A_1+\cdots+x_nA_n]^{m_1+\cdots+m_n}} .
  \label{eq:6.42}
\end{equation}
这个公式甚至在 $m_i$ 不是整数时也成立；我们将在第~10.5 节在这样一种情形下应用它。

\subsection{形状因子的计算}

现在让我们把公式 \eqref{eq:6.41} 应用于 \eqref{eq:6.38} 的分母：
\begin{equation}
  \frac{1}{((k-p)^2+i\epsilon)(k'^2-m^2+i\epsilon)(k^2-m^2+i\epsilon)}
  = \int dx\,dy\,dz\ \delta(x+y+z-1)\ \frac{2}{D^3} ,
\end{equation}
其中新的分母 $D$ 为
\begin{equation}
  \begin{aligned}
    D &= x(k^2-m^2) + y(k'^2-m^2) + z(k-p)^2 + (x+y+z)i\epsilon\\
    &= k^2 + 2k\cdot(yq-zp) + yq^2 + zp^2 - (x+y)m^2 + i\epsilon .
  \end{aligned}
\end{equation}
在第二行我们使用了 $x+y+z=1$ 和 $k' = k+q$。现在平移 $k$ 以配平方：
\begin{equation}
  \ell = k + yq - zp .
\end{equation}
经过一点代数运算，我们发现 $D$ 化简为
\begin{equation}
  D = \ell^2 - \Delta + i\epsilon,
  \label{eq:6.43}
\end{equation}
其中
\begin{equation}
  \Delta = -xyq^2 + (1-z)^2m^2 .
  \label{eq:6.44}
\end{equation}
由于对散射过程 $q^2 < 0$，$\Delta$ 是正的；我们可以把它看作一个有效质量项。

接下来我们必须用 $\ell$ 表示 \eqref{eq:6.38} 的分子。注意到由于 $D$ 只依赖于 $\ell$ 的大小，这一任务得以简化，
\begin{equation}
  \int \frac{d^4\ell}{(2\pi)^4}\ \frac{\ell^\mu}{D^3} = 0;
  \label{eq:6.45}
\end{equation}
\begin{equation}
  \int \frac{d^4\ell}{(2\pi)^4}\ \frac{\ell^\mu\ell^\nu}{D^3}
    = \frac{g^{\mu\nu}}{4}\int \frac{d^4\ell}{(2\pi)^4}\ \frac{\ell^2}{D^3} .
  \label{eq:6.46}
\end{equation}
第一个恒等式由对称性得出。为证明第二个，注意到除非 $\mu = \nu$，否则积分由对称性为零。因此 Lorentz 不变性要求我们得到正比于 $g^{\mu\nu}$ 的量。为检验系数，把两边与 $g_{\mu\nu}$ 缩并。利用这些恒等式，我们有
\begin{equation}
  \begin{aligned}
    \text{Numerator} &= \bar{u}(p')\Bigl[ \gamma^\mu\ell^2 + m^2\gamma^\mu
      - 2m(k+k')^\mu \Bigr]u(p)\\
    &\to \bar{u}(p')\Bigl[ -\tfrac{1}{2}\gamma^\mu\ell^2
      + \cdots \Bigr]u(p),
  \end{aligned}
\end{equation}
（记住 $k' = k+q$。）

把分子整理成有用的形式现在只是一些冗长的 Dirac 代数（大约一两页）的事情。这正是我们在上一节的工作得到回报之处，因为它告诉我们应该预期什么类型的答案。我们最终想把一切归并为两个分别正比于 $\gamma^\mu$ 和 $i\sigma^{\mu\nu}q_\nu$ 的项。实现这一目标最直接的方式是转而瞄准如下形式的表达式
\begin{equation}
  \gamma^\mu A + (p'^{\mu}+p^\mu)B + q^\mu C,
\end{equation}
正如在式 (6.31) 中那样。达到这个形式只需要反对易关系（例如，$\slashed{p}\gamma^\mu = 2p^\mu - \gamma^\mu\slashed{p}$）和 Dirac 方程 $\slashed{p}u(p) = mu(p)$ 与 $\bar{u}(p')\slashed{p}' = \bar{u}(p')m$；注意这蕴含 $\bar{u}(p')\slashed{q}u(p) = 0$。记住 $x+y+z = 1$ 也很有用。当尘埃落定，我们有
\begin{equation}
  \begin{aligned}
    \text{Numerator} &= \bar{u}(p')\Bigl[ \gamma^\mu\,
      \bigl( \tfrac{1}{2}\ell^2 + (1-x)(1-y)q^2 + (1-2z-z^2)m^2 \bigr)\\
    &\quad + (p'^{\mu}+p^\mu)\cdot mz(z-1)
      + q^\mu\cdot m(z-2)(x-y) \Bigr]u(p) .
  \end{aligned}
\end{equation}
根据 Ward 恒等式，$q^\mu$ 的系数必须为零，正如式 (6.31) 之后所讨论的。为看到确实如此，从 \eqref{eq:6.44} 注意分母在 $x\leftrightarrow y$ 下是对称的。$q^\mu$ 的系数在 $x\leftrightarrow y$ 下是奇性的，因此在 $x$ 和 $y$ 上积分时为零。

仍然遵循我们在上一节的工作，现在我们用 Gordon 恒等式 \eqref{eq:6.32} 消去 $(p'+p)$，代之以 $i\sigma^{\mu\nu}q_\nu$。于是我们关于电子顶角 $O(\alpha)$ 贡献的整个表达式变为
\begin{equation}
  \begin{aligned}
    \bar{u}(p')\,\delta\Gamma^\mu(p',p)\,u(p)
    &= 2ie^2 \int dx\,dy\,dz\ \delta(x+y+z-1)\ \int
      \frac{d^4\ell}{(2\pi)^4}\ \frac{1}{D^3}\\
    &\quad \times \bar{u}(p')\Bigl[
      \gamma^\mu\Bigl( \tfrac{1}{2}\ell^2 + (1-x)(1-y)q^2
        + (1-4z+z^2)m^2 \Bigr)\\
      &\quad + \frac{i\sigma^{\mu\nu}q_\nu}{2m}\,
        \bigl( 2m^2z(1-z) \bigr) \Bigr]u(p),
  \end{aligned}
  \label{eq:6.47}
\end{equation}
其中如前，
\begin{equation}
  D = \ell^2 - \Delta + i\epsilon,\qquad
  \Delta = -xyq^2 + (1-z)^2m^2 > 0 .
\end{equation}
分解成形状因子现在已一目了然。

大部分工作已经完成，我们剩下的主要任务是完成动量积分。把 $\ell^0$ 积分作为围道积分来求值，然后通过威克转动完成剩下的三维积分，这并不困难。$\ell^0$ 围道的转动如图~\ref{fig:6.1} 所示。

\begin{figure}[t]
  \centering
  \includegraphics[width=0.55\textwidth]{fig6.1.png}
  \caption{$\ell^0$ 积分的围道可以如图所示转动。}
  \label{fig:6.1}
\end{figure}

动量积分的结果为
\begin{equation}
  \int \frac{d^4\ell}{(2\pi)^4}\ \frac{1}{(\ell^2-\Delta+i\epsilon)^3}
    = \frac{-i}{32\pi^2\Delta},
  \qquad
  \int \frac{d^4\ell}{(2\pi)^4}\ \frac{\ell^2}{(\ell^2-\Delta+i\epsilon)^3}
    = \frac{i}{16\pi^2}\ \cdots ,
  \label{eq:6.48}
\end{equation}
将其代入 \eqref{eq:6.47}，即得一圈形状因子。在紫外区域，积分对数发散；我们用截断 $\Lambda$（或维数正则化）来正规化它们。在这两种临时修正之下，形状因子为
\begin{equation}
  F_1(q^2) = 1 + \frac{\alpha}{2\pi}\int dx\,dy\,dz\
    \delta(x+y+z-1)\ \log\Bigl(
      \frac{m^2(1-z)^2}{m^2(1-z)^2 - q^2xy} \Bigr)\
    + \frac{\alpha}{2\pi}\int dx\,dy\,dz\ \delta(x+y+z-1)\
      \frac{m^2(1-4z+z^2) + q^2(1-x)(1-y)}
        {m^2(1-z)^2 - q^2xy + \mu^2 z}\
    + \mathcal{O}(\alpha^2);
  \label{eq:6.56}
\end{equation}
\begin{equation}
  F_2(q^2) = \frac{\alpha}{2\pi}\int dx\,dy\,dz\
    \delta(x+y+z-1)\ \frac{2m^2z(1-z)}{m^2(1-z)^2 - q^2xy}
    + \mathcal{O}(\alpha^2) .
  \label{eq:6.57}
\end{equation}
注意紫外和红外发散都不影响 $F_2(q^2)$。因此我们可以毫不含糊地计算
\begin{equation}
  F_2(q^2 = 0) = \frac{\alpha}{2\pi}\int_0^1 dz\ \int_0^{1-z} dy\
    \frac{2m^2z(1-z)}{m^2(1-z)^2}
    = \frac{\alpha}{2\pi}\ \int_0^1 dz\ \frac{z}{1-z}\ \cdots
    = \frac{\alpha}{2\pi} .
  \label{eq:6.58}
\end{equation}
于是，我们得到电子 $g$ 因子的一个修正：
\begin{equation}
  \frac{g-2}{2} = \frac{\alpha}{2\pi} = 0.0011614 .
  \label{eq:6.59}
\end{equation}
这个结果最早由 Schwinger 于 1948 年得到。%
\footnote{J.~Schwinger, \emph{Phys.\ Rev.} \textbf{73}, 416L (1948)。}
实验给出 $a_e = 0.0011597$。显然，我们为 $F_2(0)$ 得到的这个不含糊的值，直到 $\alpha$ 的高阶，也是不含糊地正确的。

\subsection{QED的精密检验}

在 $\alpha$ 阶 QED 对 $a_e$ 的预言取得成功的基础上，一代又一代物理学家提高了对该量理论和实验确定的精度。$a_e$ 的 QED 公式系数现在已知到 $\alpha^4$ 阶。$\alpha^2$ 阶及更高阶系数的计算需要对紫外发散作系统处理。

这些富有挑战性的理论计算伴随着越来越富有想象力的实验。最近一次对 $a_e$ 的测量使用了一种由 Dehmelt 及其合作者发展的技术，其中单个电子被囚禁在静电场和静磁场系统中，并被激发到自旋共振。%
\footnote{R.~Van Dyck, Jr., P.~Schwinberg, and H.~Dehmelt, \emph{Phys.\ Rev.\
Lett.} \textbf{59}, 26 (1987)。} 如今，$a_e$ 的最佳理论和实验值一致到八位有效数字。

高阶 QED 计算也已针对其他若干量完成。这些包括氢原子和类氢原子的跃迁能量、缪子的反常磁矩，以及单态和三重态电子偶素的衰变率。这些量中有许多也已得到高精度测量。这些比较的全部集合在各种情形下对 QED 的有效性给出详细的检验。这些精密检验的结果汇总在表~\ref{tab:6.1} 中。

报告 QED 理论与实验之间精密比较的结果有一些微妙之处，因为理论预言需要一个极其精确的 $\alpha$ 值，而这个值只能从另一个精密 QED 实验获得。因此，我们把理论与实验之间的每次比较都引用为对 $\alpha$ 的一次独立确定。每个 $\alpha$ 值被赋予一个误差，它是来自理论和实验的预期不确定性的复合。QED 在多大程度上得到证实，取决于来自不同来源的 $\alpha$ 值在多大程度上一致。

表~\ref{tab:6.1} 中前九项涉及原子物理情境中的 QED 计算。其中，利用 Ramsey 的氢微波激射器测量的氢原子超精细劈裂，是物理学中已知最精确的量。遗憾的是，质子内部结构的影响带来不确定性，限制了理论上预测这个量的精度。同样的困难也适用于兰姆位移，即氢原子 $j = 1/2$ 的 $2S$ 和 $2P$ 能级之间的劈裂。现在最精确的 QED 检验来自不涉及强相互作用粒子的系统，即电子 $g-2$ 和 $e^{-}\mu^{+}$ 原子（缪子素）中的超精细劈裂。这一组中的最后一项给出一种确定 $\alpha$ 的新方法，它利用精确已知的质量比，把对中子康普顿波长的一个非常精确的测量换算成电子质量的值。这可以与已知的里德伯能量值和精确的 QED 公式结合来确定 $\alpha$。这些数值之间唯一严重的差异出现在三重态电子偶素衰变率上；然而，有一些证据表明，相对阶 $\alpha^2$ 的图对表中所引用的值给出大的修正。

接下来两项是从高能电子对撞机上高阶 QED 反应对 $\alpha$ 的确定。这些高能实验通常只能达到百分之一的精度，但它们的结论与较低能量下可获得的精确信息一致。

最后，表中最后两项给出对 $\alpha$ 的两个独立测量，它们来自凝聚态系统中奇异的量子干涉现象。这两个效应分别提供一个标准电阻和一个标准频率，人们相信它们以对宏观系统严格为零的修正来测量电子的电荷。%
\footnote{关于这些效应及其与 $\alpha$ 的精确关系，参见 D.~R.~Yennie, \emph{Rev.\ Mod.\ Phys.} \textbf{59}, 781 (1987)。}

\begin{table}[t]
  \centering
  \caption{由精密 QED 实验得到的 $1/\alpha$ 值}
  \begin{tabular}{ll}
    \toprule
    低能 QED： & \\
    电子 $(g-2)$ & 137.035\,992\,35\,(73)\\
    缪子 $(g-2)$ & 137.035\,5\,(11)\\
    缪子素超精细劈裂 & 137.035\,994\,(18)\\
    兰姆位移 & 137.036\,8\,(7)\\
    氢原子超精细劈裂 & 137.036\,0\,(3)\\
    电子偶素中的 $2\,^3S_1 - 1\,^3S_1$ 劈裂 & 137.034\,(16)\\
    $1\,S_0$ 电子偶素衰变率 & 137.00\,(6)\\
    $3\,S_1$ 电子偶素衰变率 & 136.971\,(6)\\
    中子康普顿波长 & 137.036\,010\,1\,(54)\\
    \midrule
    高能 QED： & \\
    $\sigma(e^{+}e^{-}\to e^{+}e^{-}e^{+}e^{-})$ & 136.5\,(2.7)\\
    $\sigma(e^{+}e^{-}\to e^{+}e^{-}\mu^{+}\mu^{-})$ & 139.9\,(1.2)\\
    \midrule
    凝聚态物质： & \\
    量子霍尔效应 & 137.035\,997\,9\,(32)\\
    交流约瑟夫森效应 & 137.035\,977\,0\,(77)\\
    \bottomrule
  \end{tabular}
  \label{tab:6.1}
\end{table}

表中显示的每个 $\alpha$ 值都是通过把实验测量拟合到含参数 $\alpha$ 的理论表达式而得到的。括号中的数字是最后显示数字位的标准误差，包括理论和实验的不确定性。本表基于 Kinoshita (1990) 的精密 QED 综述中给出的结果。该书包含一系列清晰综述，介绍为详细分析 QED 过程而发展起来的卓越理论与实验技术。五个最精确的值按照 T.~Kinoshita 在 \emph{History of Original Ideas and Basic Discoveries in Particle Physics}（H.~Newman 和 T.~Ypsilantis 主编，Plenum Press, New York, 1995）中的给出做了更新。后一篇论文还对精密 QED 实验的未来给出了有趣的展望。

整个图景契合得远远超出任何合理预期。基于此表给出的证据，QED 是所有物理理论中被检验得最严格——也是最惊人地成功——的。

\section{电子顶角函数：红外发散}

现在让我们面对 $F_1(q^2)$ 结果 \eqref{eq:6.56} 中的红外发散。在 $\mu\to 0$ 极限下的主导部分是
\begin{equation}
  F_1(q^2) \sim \frac{\alpha}{2\pi}\int dx\,dy\,dz\ \delta(x+y+z-1)\
    \frac{m^2(1-4z+z^2) + q^2(1-x)(1-y)}
      {m^2(1-z)^2 - q^2xy + \mu^2z}
    + \cdots .
  \label{eq:6.60}
\end{equation}
为理解这个表达式，我们必须做一些化简工作，提取并求值积分的发散部分。在本节中，我们只保留在 $\mu\to 0$ 极限下发散的项。

首先注意发散发生在费曼参数空间 $z\approx 1$（因此 $x\approx y\approx 0$）的角落。在这个区域，我们可以在 \eqref{eq:6.60} 的分子中取 $z = 1$ 和 $x = y = 0$。我们还可以在分母的 $\mu^2$ 项中取 $z = 1$。利用 delta 函数求 $x$ 积分，我们于是有
\begin{equation}
  F_1(q^2) = \frac{\alpha}{2\pi}\int_0^1 dz\ \int_0^{1-z} dy\
    \Bigl[ \frac{-2m^2 + q^2}{m^2(1-z)^2 - q^2y(1-z-y) + \mu^2}
      + \frac{m^2}{(1-z)^2 + \mu^2/m^2} \Bigr].
\end{equation}
（$z$ 积分的下限不重要。）作变量替换
\begin{equation}
  v = 1 - z,\qquad w = \frac{y}{1-z},
\end{equation}
这个表达式变为
\begin{equation}
  F_1(q^2) \sim \frac{\alpha}{2\pi}\int_0^1 dv\
    \int_0^1 dw\ \Bigl[ \frac{-2m^2}{m^2w^2 + \mu^2}
      + \frac{m^2}{m^2v^2 + \mu^2}\ \cdots \Bigr],
\end{equation}
\begin{equation}
  \sim \frac{\alpha}{2\pi}\ \Bigl[ \cdots \Bigr]
    \log\Bigl( \frac{-q^2}{\mu^2} \Bigr)\log\Bigl( \frac{-q^2}{m^2}
      \Bigr) + \cdots ,
  \label{eq:6.61}
\end{equation}
这就是第~6.1 节中发现的苏达科夫双对数。这个发散必须与轫致辐射贡献相抵消的物理要求，决定了结果的减除方式和物理解释。

\section{红外发散求和与解释}

红外发散来自具有``软''动量的光子：能量小于某个截断 $E_{-}$ 的实光子，以及（在威克转动后）$k^2 < E_{-}^2$ 的虚光子。一个典型的高阶图将涉及大量的实光子和虚光子。但要找到一个发散，我们需要的不仅仅是一个软光子；我们需要电子传播子中一个奇异的分母。例如，考虑以下两个图：
\begin{equation}
  \text{(soft photon emitted before hard photon)} \qquad
  \text{(soft photon emitted after hard photon)} .
\end{equation}
第一个图（电子先发射一个软光子，随后发射一个硬光子）没有红外发散，因为两个电子传播子中的动量都远离质壳。然而，如果软光子最后发射，相邻传播子的分母是 $(p'+k)^2 - m^2 = 2p'\cdot k$，它随 $k\to 0$ 趋于零。因此第二个图确实包含一个发散。于是，我们想要考虑这样的图：一个任意的硬过程（可能涉及硬光子和软光子的发射），被在电子腿上添加软的实光子和虚光子所修改：
\begin{equation}
  \underbrace{\text{soft real and virtual photons}}_{\text{on the electron
    legs}}\quad \underbrace{\text{arbitrary hard process}} .
\end{equation}
遵循 Weinberg，我们将把所有这类图的贡献加起来。这个计算中唯一新的困难在于对光子可以出现的所有方式进行计数的组合问题。

首先考虑出射电子线。我们在该线上连接 $n$ 个光子，动量为 $k_1, \ldots, k_n$。目前我们不关心这些是外光子、彼此相连的虚光子，还是连接到入射电子线上顶点的虚光子。这个图的 Dirac 结构为
\begin{equation}
  \bar{u}(p')\,(-ie\gamma^{\mu_1})\,
  \frac{i(\slashed{p}'+\slashed{k}_1+m)}{2p'\cdot k_1 + \mathcal{O}(k^2)}\,
  (-ie\gamma^{\mu_2})\,
  \frac{i(\slashed{p}'+\slashed{k}_1+\slashed{k}_2+m)}
    {2p'\cdot(k_1+k_2) + \mathcal{O}(k^2)}\ \cdots\,
  (-ie\gamma^{\mu_n})\,
  \frac{i(\slashed{p}'+\textstyle\sum\slashed{k}_i+m)}
    {2p'\cdot(k_1+\cdots+k_n) + \mathcal{O}(k^2)}\,
  \mathcal{M}_{\mathrm{hard}} .
  \label{eq:6.73}
\end{equation}
我们将假设所有 $k_i$ 都小，丢弃分母中的 $\mathcal{O}(k^2)$ 项。我们还将丢弃分子中的 $\slashed{k}$ 项，正如我们在第~6.1 节处理轫致辐射时那样。此外，正如我们在那里所做的，我们可以把因子 $(\slashed{p}'+m)$ 推到左边，并利用 $\bar{u}(p')(\slashed{p}'-m) = 0$：
\begin{equation}
  \bar{u}(p')\slashed{\epsilon}^{\mu_1}(\slashed{p}'+m)
    \slashed{\epsilon}^{\mu_2}(\slashed{p}'+m)\cdots
  = \bar{u}(p')\,2p'\cdot\epsilon^{\mu_1}\slashed{\epsilon}^{\mu_2}
    (\slashed{p}'+m) = \bar{u}(p')\,2p'\cdot\epsilon^{\mu_1}\,
    2p'\cdot\epsilon^{\mu_2}\cdots .
\end{equation}
这把表达式 \eqref{eq:6.73} 变为
\begin{equation}
  \bar{u}(p')\,(-ie)^n\ \prod_{i=1}^{n}\ \frac{p'\cdot\epsilon^{\mu_i}}
    {p'\cdot(k_1+\cdots+k_i)}\ \cdots .
  \label{eq:6.74}
\end{equation}
仍然只处理出射电子线，我们现在必须对动量 $k_1\cdots k_n$ 的所有可能排列求和。（当其中两个光子连在一起形成一个虚光子时，这一过程会重复计数。我们稍后处理这种重复计数。）共有 $n!$ 个不同的图需要求和，对应于 $n$ 个光子动量的 $n!$ 个排列。设 $\pi$ 表示这样一个排列，于是 $\pi(i)$ 是 $i$ 被映到的 1 与 $n$ 之间的数。（例如，如果 $\pi$ 表示把 $1\to 3$、$2\to 1$ 和 $3\to 2$ 的排列，则 $\pi(1) = 3$、$\pi(2) = 1$、$\pi(3) = 2$。）

有了这个记号，我们可以借助以下恒等式完成对排列的求和：
\begin{equation}
  \sum_{\text{all permutations }\pi}\ \frac{1}{p\cdot k_{\pi(1)}}\,
    \frac{1}{p\cdot(k_{\pi(1)}+k_{\pi(2)})}\ \cdots\,
    \frac{1}{p\cdot(k_{\pi(1)}+\cdots+k_{\pi(n)})}
  = \frac{1}{p\cdot k_1\,p\cdot k_2\cdots p\cdot k_n} .
  \label{eq:6.75}
\end{equation}
这个公式的证明对 $n$ 作归纳。对 $n=2$ 我们有
\begin{equation}
  \frac{1}{p\cdot k_{\pi(1)}}\frac{1}{p\cdot(k_{\pi(1)}+k_{\pi(2)})}
  = \frac{1}{p\cdot k_1\,p\cdot k_2},
\end{equation}
这通过把两项相加很容易验证。对于归纳步骤，注意 \eqref{eq:6.75} 左边的最后一个因子对每个排列 $\pi$ 都相同。把这个因子提到求和号外，左边变为
\begin{equation}
  \text{LHS} = \sum_{\pi}\ \frac{1}{p\cdot k_{\pi(1)}}\ \cdots
    \frac{1}{p\cdot(k_{\pi(1)}+\cdots+k_{\pi(n-1)})}\ .
\end{equation}
对任意给定的 $\pi$，被求和的量都与 $k_{\pi(n)}$ 无关。令 $i = \pi(n)$，我们现在可以写成
\begin{equation}
  \sum_{\pi} = \sum_{i=1}^{n}\ \sum_{\pi'(i)},
\end{equation}
其中 $\pi'(i)$ 是其余 $n-1$ 个整数上所有排列的集合。由归纳假设 \eqref{eq:6.75} 对 $n-1$ 成立，我们有
\begin{equation}
  \text{LHS} = \frac{1}{p\cdot k_1\,p\cdot k_2\cdots p\cdot k_{i-1}\,
    p\cdot k_{i+1}\cdots p\cdot k_n}\ .
\end{equation}
如果我们现在把此和中每一项都乘以再除以 $p\cdot k_i$，就很容易得到我们想要的结果 \eqref{eq:6.75}。

把 \eqref{eq:6.75} 应用于 \eqref{eq:6.74}，我们发现
\begin{equation}
  \bar{u}(p')\,(-ie)^n\ \Bigl( \frac{p'\cdot\epsilon_1}{p'\cdot k_1}
    \Bigr)\Bigl( \frac{p'\cdot\epsilon_2}{p'\cdot k_2} \Bigr)\cdots
    \Bigl( \frac{p'\cdot\epsilon_n}{p'\cdot k_n} \Bigr)
  = \bar{u}(p')\prod_{i=1}^{n}\ \Bigl( \frac{-ie\,p'\cdot\epsilon_i}
    {p'\cdot k_i} \Bigr),
  \label{eq:6.76}
\end{equation}
其中圆圈表示对插入 $n$ 条光子线的所有可能顺序求和。

一组类似的运算化简了在初始电子线上软光子插入的求和。然而在那里，传播子动量为 $p-k_1$、$p-k_1-k_2$，等等。因此每个光子的因子中我们得到一个额外的负号，因为
\begin{equation}
  (p-\textstyle\sum k)^2 - m^2 \approx -2p\cdot\textstyle\sum k .
\end{equation}
现在考虑总共包含 $n$ 个软光子的图，它们以任何可能顺序连接到初始或最终电子线。所有这类图的和可以写成
\begin{equation}
  \bar{u}(p')\,\mathcal{M}_{\mathrm{hard}}\,u(p)\ \prod_{i=1}^{n}\ (-ie)\
    \Bigl( \frac{p'\cdot\epsilon_i}{p'\cdot k_i} - \frac{p\cdot\epsilon_i}
      {p\cdot k_i} \Bigr),
  \label{eq:6.77}
\end{equation}
通过把所有因子乘开，你可以看到我们为把 $n$ 个光子分配给两条线的每种可能方式得到正确的项。

接下来我们必须确定哪些光子是实的、哪些是虚的。我们可以通过选取两个光子动量 $k_i$ 和 $k_j$、令 $k_j = -k_i = k$、乘以光子传播子并对 $k$ 积分，来产生一个虚光子。对每个虚光子我们于是得到表达式
\begin{equation}
  \begin{aligned}
    X &= \frac{1}{2}\ \frac{e^2}{2}\ \int \frac{d^4k}{(2\pi)^4}\
      \frac{-i}{k^2+i\epsilon}\ \Bigl( \frac{p'}{p'\cdot k}
        - \frac{p}{p\cdot k} \Bigr)\Bigl( \frac{p'}{-p'\cdot k}
          - \frac{p}{-p\cdot k} \Bigr)\\
    &= \frac{ie^2}{2}\ \int \frac{d^4k}{(2\pi)^4}\ \frac{1}{k^2+i\epsilon}\
      \Bigl( \frac{p'}{p'\cdot k} - \frac{p}{p\cdot k} \Bigr)^2 .
  \end{aligned}
  \label{eq:6.78}
\end{equation}
因子 $1/2$ 是必需的，因为我们的程序把每个费曼图计数了两次：交换 $k_i$ 和 $k_j$ 得到同一个图。通过仔细的围道积分可以求值这个表达式，但有一种更简单的方式。注意这个近似方案给具有一个圈且无外光子的图赋予值
\begin{equation}
  \bar{u}(p')(i\mathcal{M}_{\mathrm{hard}})u(p)\cdot X .
\end{equation}
因此，$X$ 必定正是 \eqref{eq:6.61} 中所显示的形状因子一圈修正的红外极限：
\begin{equation}
  X = \frac{\alpha}{\pi}\mathcal{I}_{\mathrm{IR}}(q^2)\log\Bigl(
    \frac{-q^2}{\mu^2} \Bigr) .
  \label{eq:6.79}
\end{equation}
从 \eqref{eq:6.78} 直接推导这个结果见上文引用的 Weinberg 的论文。注意，在我们上一节的论证中，结果 \eqref{eq:6.79} 只是在 $q^2 = 0$ 处减除之后才得出的，因此我们应该担心 \eqref{eq:6.79} 是否与第 $n$ 阶图的相应减除相容。此外，我们正在求和的某些图包含我们没有讨论过的外腿修正。这里我们只是说明，这两种微妙之处都不影响最终答案；其证明需要 Yennie、Frautschi 和 Suura 论文中的重工具。

如果有 $m$ 个虚光子，我们得到 $m$ 个像 \eqref{eq:6.79} 的因子，此外还有一个 $1/m!$ 的对称因子，因为虚光子彼此交换不改变图。于是我们可以对 $m$ 求和，得到由于任意多个软虚光子的存在而产生的完整修正：
\begin{equation}
  \sum_{m=0}^{\infty}\ \frac{X^m}{m!}
    = \bar{u}(p')(i\mathcal{M}_{\mathrm{hard}})u(p)\exp(X) .
  \label{eq:6.80}
\end{equation}
如果在 $m$ 个虚光子之外我们还发射一个实光子，我们必须乘以其极化矢量、对极化求和，并把矩阵元的平方对光子的相空间积分。这给出一个额外因子
\begin{equation}
  \int \frac{d^3k}{(2\pi)^3}\ \frac{1}{2k}\ \frac{\alpha}{2\pi}\ \sum
    \Bigl| \frac{p'\cdot\epsilon}{p'\cdot k} - \frac{p\cdot\epsilon}{p\cdot k}
      \Bigr|^2
  = \frac{\alpha}{\pi}\ \mathcal{I}(\mathbf{v},\mathbf{v}')\log\Bigl(
    \frac{E_{-}}{\mu} \Bigr)
  = Y ,
  \label{eq:6.81}
\end{equation}
在截面中。假设光子的能量大于 $\mu$ 且小于 $E_{-}$（探测器阈值），这个表达式就是
\begin{equation}
  Y = \frac{\alpha}{\pi}\mathcal{I}(\mathbf{v},\mathbf{v}')\log\Bigl(
    \frac{E_{-}}{\mu} \Bigr)
  = \frac{\alpha}{\pi}\mathcal{I}_{\mathrm{IR}}(q^2)\log\Bigl(
    \frac{E_{-}}{\mu} \Bigr) .
  \label{eq:6.82}
\end{equation}
如果发射 $n$ 个实光子，我们得到 $n$ 个这样的因子，还有一个 $1/n!$ 的对称因子，因为末态中有 $n$ 个全同玻色子。因此发射任意数量软光子的截面为
\begin{equation}
  \sum_{n=0}^{\infty}\ \frac{Y^n}{n!}\ d\sigma_{\mathrm{hard}}
    = d\sigma_{\mathrm{hard}}\ \exp(Y) .
  \label{eq:6.83}
\end{equation}
不发射硬光子而散射的概率就是总截面，而表达式
\begin{equation}
  d\sigma_{\mathrm{total}} = d\sigma_{\mathrm{hard}}\,
    \exp(X + Y) ,
  \label{eq:6.84}
\end{equation}
是本节的核心结果。$X$（来自虚光子）和 $Y$（来自实光子）的红外发散精确抵消，留下一个有限的结果。在大动量转移极限下，指数因子包含苏达科夫双对数，且
\begin{equation}
  \exp(X+Y) \sim \exp\Bigl[ -\frac{\alpha}{\pi}
    \mathcal{I}_{\mathrm{IR}}(q^2)\log\Bigl( \frac{-q^2}{m^2} \Bigr)
      \log\Bigl( \frac{-q^2}{E_{-}^2} \Bigr) \Bigr] .
  \label{eq:6.85}
\end{equation}
在这个极限下，不发射硬光子而散射的概率比 $q^2$ 的任何幂次都更快地减小。这个包含苏达科夫双对数的指数修正因子称为\emph{苏达科夫形状因子}。

为结束本节，让我们在同样的近似下计算某个硬散射过程伴随着产生 $n$ 个能量都在 $E_{-}$ 与 $E_{+}$ 之间的软光子的概率。这些光子的相空间积分给出 $\log(E_{+}/E_{-})$ 而不是 $\log(E_{\mathrm{max}}/\mu)$。如果我们把能量大于 $E_{+}$ 的光子划归过程的``硬''部分，我们发现截面由 \eqref{eq:6.84} 乘以额外因子给出
\begin{equation}
  \mathrm{Prob}\bigl( n\gamma \text{ with } E_{-} < E < E_{+} \bigr)
    = \frac{1}{n!}\Bigl[ \frac{\alpha}{\pi}\mathcal{I}_{\mathrm{IR}}(q^2)
      \log\Bigl( \frac{E_{+}}{E_{-}} \Bigr) \Bigr]^n
      \exp\Bigl[ -\frac{\alpha}{\pi}\mathcal{I}_{\mathrm{IR}}(q^2)
      \log\Bigl( \frac{E_{+}}{E_{-}} \Bigr) \Bigr] .
  \label{eq:6.86}
\end{equation}
这个表达式具有泊松分布的形式，
\begin{equation}
  P(n) = \frac{1}{n!}\lambda^n e^{-\lambda},
\end{equation}
其中
\begin{equation}
  \lambda = \frac{\alpha}{\pi}\mathcal{I}_{\mathrm{IR}}(q^2)\log\Bigl(
    \frac{E_{+}}{E_{-}} \Bigr) .
\end{equation}
这正是我们在式~\eqref{eq:6.19} 中对辐射光子数所做的半经典估计。

\problems

\textbf{6.1} \emph{Rosenbluth 公式。} 如第~6.2 节所讨论，Dirac 费米子的精确电磁相互作用顶角可以相当一般地用两个形状因子 $F_1(q^2)$ 和 $F_2(q^2)$ 写成：
\begin{equation}
  -ie\Gamma^\mu(p',p) = -ie\Bigl[ \gamma^\mu F_1(q^2)
    + \frac{i\sigma^{\mu\nu}q_\nu}{2m}F_2(q^2) \Bigr],
\end{equation}
其中 $q = p'-p$，$\sigma^{\mu\nu} = \tfrac{i}{2}[\gamma^\mu, \gamma^\nu]$。如果费米子是一个强相互作用粒子，例如质子，则形状因子反映由强相互作用产生的结构，因此不容易从第一性原理计算。然而，这些形状因子可以通过实验确定。考虑一个能量 $E \gg m_e$ 的电子在最初静止的质子上的散射。证明上述顶角表达式导致如下弹性散射截面表达式（\emph{Rosenbluth 公式}），它计算到 $\alpha$ 的最低阶，但对强相互作用则到所有阶：
\begin{equation}
  \frac{d\sigma}{d\Omega} = \frac{\alpha^2\cos^2(\theta/2)}
    {4E^2\sin^4(\theta/2)}\ \frac{1}{1 + (2E/m_p)\sin^2(\theta/2)}\
    \Bigl[ F_1^2 - \frac{q^2}{4m_p^2}\ \Bigl( F_2^2 \Bigr)
      + \cdots \Bigr],
\end{equation}
其中 $\theta$ 是实验室系散射角，$F_1$ 和 $F_2$ 在此角度弹性散射所对应的 $q^2$ 处取值。通过测量 $(d\sigma/d\cos\theta)$ 随角度的变化，就有可能提取 $F_1$ 和 $F_2$。注意当 $F_1 = 1$ 且 $F_2 = 0$ 时，Rosenbluth 公式退化为点粒子散射的 Mott 公式（无质量极限下）（参见习题~5.1）。

\textbf{6.2} \emph{等效光子近似。} 考虑非常高能的电子在靶上散射的过程。在 $\alpha$ 的最低阶，电子通过一个光子传播子连接到靶。如果电子的初末能量分别为 $E$ 和 $E'$，光子将携带动量 $q$，使得 $q^2 \approx -2EE'(1-\cos\theta)$。在前向散射极限下，无论能量损失如何，光子动量都趋近 $q^2 = 0$；因此反应在前向方向高度成峰。人们忍不住猜想，在这个极限下虚光子变成实光子。让我们
\begin{enumerate}[label=(\alph*)]
  \item 证明：对于电子在靶上散射（假设靶无自旋且非常重）的过程，微分截面可以写成光子-靶截面与一个电子-光子通量因子的乘积。
  \item 计算前向方向的虚光子通量，并证明它具有等效光子谱的形式。
  \item 证明 Weizs\"acker--Williams 近似在光子几乎为实的极限下对总截面给出很好的描述。
  \item 推导虚光子能谱的 Weizs\"acker--Williams 公式。
\end{enumerate}

\textbf{6.3} \emph{对 $g-2$ 的奇异贡献。} 任何与电子耦合的粒子都会对电子-光子形状因子产生修正，特别是对 $g-2$ 产生修正。由于电子 $g-2$ 与 QED 高精度一致，这些修正使我们能够约束假设的新粒子的性质。
\begin{enumerate}[label=(\alph*)]
  \item 弱相互作用和电磁相互作用的统一理论包含一个称为 Higgs 玻色子的标量粒子 $h$，它按照下式与电子耦合
    \begin{equation}
      \mathcal{L} = \lambda\bar{\psi}\psi h .
    \end{equation}
    用 $\lambda$ 和 Higgs 玻色子的质量 $m_h$ 计算一个虚 Higgs 玻色子对电子 $(g-2)$ 的贡献。
  \item QED 对电子的反常磁矩解释得极好。如果 $a = (g-2)/2$，
    \begin{equation}
      |a_{\mathrm{expt.}} - a_{\mathrm{QED}}| < 1\times10^{-10} .
    \end{equation}
    这给 $\lambda$ 和 $m_h$ 施加了什么限制？在最简单的电弱理论版本中，$\lambda = 3\times10^{-6}$ 且 $m_h > 60$ GeV。证明这些值没有被排除。Higgs 玻色子与缪子的耦合大一个因子 $(m_\mu/m_e)$：$\lambda = 6\times10^{-4}$。因此，尽管我们对缪子反常磁矩的实验认识不那么精确，
    \begin{equation}
      |a_{\mu,\mathrm{expt.}} - a_{\mu,\mathrm{QED}}| < 3\times10^{-9},
    \end{equation}
    仍然可以得到对 $m_h$ 更强的限制。它足够强吗？
  \item 这一理论的一些更复杂的版本包含一个称为轴子的赝标量粒子，它按照下式与电子耦合
    \begin{equation}
      \mathcal{L} = \lambda\bar{\psi}\gamma^5\psi\,a .
    \end{equation}
    轴子可能像电子一样轻，甚至更轻，并且可能比 Higgs 玻色子耦合更强。计算一个虚轴子对电子 $g-2$ 的贡献，并求出被排除的 $\lambda$ 和 $m_a$ 值。
\end{enumerate}
